% Case study write-up: Berlin Hausverwaltung leads for Heidi Systems.
% Compile with pdflatex (e.g. on overleaf.com: New Project > Upload Project > this file).
\documentclass[11pt,a4paper]{article}

\usepackage[T1]{fontenc}
\usepackage[utf8]{inputenc}
\usepackage{lmodern}
\usepackage[a4paper,margin=2.2cm]{geometry}
\usepackage{booktabs}
\usepackage{tabularx}
\usepackage{array}
\usepackage[table]{xcolor}
\usepackage{enumitem}
\usepackage{titlesec}
\usepackage{fancyvrb}
\usepackage[hidelinks]{hyperref}

\newcommand{\repourl}{https://github.com/6hammad9/heidi-lead-pipeline}

\definecolor{accent}{HTML}{0B6A70}
\definecolor{assume}{HTML}{9A4F0A}
\definecolor{muted}{HTML}{56636E}
\newcommand{\ev}{\textcolor{accent}{\textsc{evidence}}}
\newcommand{\as}{\textcolor{assume}{\textsc{assumption}}}
\newcommand{\code}[1]{\texttt{#1}}

\titleformat{\section}{\large\bfseries\color{accent}}{\thesection}{0.6em}{}
\titleformat{\subsection}{\normalsize\bfseries}{\thesubsection}{0.6em}{}
\setlist{itemsep=2pt,topsep=4pt}
\setlength{\parindent}{0pt}
\setlength{\parskip}{6pt}
\renewcommand{\arraystretch}{1.25}
\newcolumntype{L}{>{\raggedright\arraybackslash}X}

\begin{document}

{\color{muted}\small Case study \textbullet{} AI Agent Builder \textbullet{} Heidi Systems}\\[2pt]
{\LARGE\bfseries Berlin Hausverwaltung Leads}\\[4pt]
{\color{muted} Sources, KPIs, automation and validation \textbullet{} 7 November 2026}

\medskip
An automated pipeline that finds Berlin property management companies, has an AI agent read
each company's website, checks every claim the agent makes against the page it came from, and
scores the companies as leads for a metering service.

\begin{center}
\begin{tabular}{cccc}
\toprule
\textbf{2,080} & \textbf{1,563} & \textbf{96\%} & \textbf{\$2.93} \\
Berlin companies & websites researched & claims confirmed on page & total AI cost \\
\bottomrule
\end{tabular}
\end{center}

\textbf{Deliverable:} \code{heidi\_leads\_berlin.xlsx} (tabs: Leads, KPI definitions, Run log).
\textbf{Code:} \url{\repourl}

%--------------------------------------------------------------------
\section{How the KPIs were chosen}

I didn't start from whatever data was easy to find. I fixed the method first, then derived the
KPIs from it, so every column in the sheet can be traced back to a question and a source.

\begin{enumerate}
  \item \textbf{Goal-Question-Metric} \cite{gqm}: one goal, five questions, one measurable KPI
        per question.
  \item \textbf{SMART check} \cite{smart}: a KPI that can't be measured or found gets demoted to
        an info column.
  \item \textbf{Three tests}: does it correlate with deal value, does it vary across companies,
        and can the data actually be found?
  \item \textbf{Weighted score}: tier weights 3\,/\,2\,/\,1, each KPI scored 0--5.
\end{enumerate}

\textbf{Goal:} find the Berlin Hausverwaltungen that are the most valuable and most winnable
customers for a metering service.

\begin{tabularx}{\linewidth}{@{}L L l@{}}
\toprule
\textbf{Question} & \textbf{KPI that answers it} & \textbf{Type} \\
\midrule
Q1 How much metering volume could they bring? & Units managed & Firmographic \\
Q2 Is their portfolio covered by HeizkostenV billing? & Residential share & Firmographic \\
Q3 Who signs the metering contract, and how easily can it be won? & Buyer type; association
  membership as a bonus & Firmographic / behavioral \\
Q4 How easily could they switch to Heidi? & Uses software Heidi integrates with & Technographic \\
Q5 Is there a reason to act now? & Not scored: the 31 Dec 2026 remote-reading deadline applies to
  the whole market (E9) & Info column \\
\bottomrule
\end{tabularx}

\subsection{The scored KPIs}

Every scoring choice is labelled \ev{} (with a source) or \as{} (not yet proven).

\begin{tabularx}{\linewidth}{@{}p{2.6cm} p{1.9cm} L L@{}}
\toprule
\textbf{KPI} & \textbf{Tier} & \textbf{Scale (0--5)} & \textbf{Why} \\
\midrule
Units managed & Critical $\times$3 & Five bands from published values: $<$800, 800--1.4k,
  1.4k--3k, 3k--6.6k, $>$6.6k & \ev{} Heidi prices per unit per year, so revenue scales with
  units (E1) \\
Residential share & Important $\times$2 & Mainly residential 5, mixed 3, mainly commercial 1 &
  \as{} Residential portfolios hold more buildings that need heating-cost billing. This is a
  proxy: the HeizkostenV also covers commercial buildings \\
Buyer type & Important $\times$2 & Rental manager 5; WEG or mixed manager 4; cooperative or own
  portfolio 3; municipal 2 & \ev{} A rental manager can sign for the owner (E4). In a WEG the
  owners vote on the provider (\S\,19 WEG, E3). Municipal companies run metering in-house
  (E6, E7, confirmed again in this run) \\
Heidi-compatible software & Useful $\times$1 & Domus, Immoware24, Wodis, Immocloud 5; other
  software 2 & \ev{} Heidi's integration list (E2); \as{} that integration raises the win rate \\
Association (VDIV, IVD, BVI) & Bonus +3 & Only adds points & \as{} A sign of a professional
  company and a possible sales channel \\
Growth signal & Bonus +3 & Only adds points & \as{} A growing portfolio means more units \\
\bottomrule
\end{tabularx}

\[
\text{Base} = \frac{\sum (\text{weight} \times \text{sub-score})}{\sum \text{weight of
\emph{known} KPIs}} \in [0,5]
\qquad
\text{Lead score} = \text{Base} \times 18.8 + \text{bonus} \in [0,100]
\]

\begin{itemize}
  \item \textbf{Unknown is not zero.} A KPI the website doesn't mention is left out of the
        average instead of counting as 0. Otherwise small companies with detailed websites would
        rank above large ones that publish little. A \emph{Data completeness} column shows how
        much of the weight is known.
  \item \textbf{Confidence flag.} Confidence is \emph{Low} when less than half of the weight is
        known or the company size is unknown.
  \item \textbf{No double counting.} Units, buildings and employees all measure size, so only
        units are scored. Employees and buildings are used to \emph{estimate} units when a company
        doesn't publish them. The ratio comes from our own data: 150 units per employee
        ($n=64$) and 29 units per building ($n=38$). Estimates are marked \emph{estimated} in
        the sheet.
  \item \textbf{Lead quality and reachability are kept apart.} Whether we found a contact person
        says nothing about how good a customer the company would be. \emph{Contact status} is
        its own column and doesn't affect the score.
\end{itemize}

\subsection{Sources behind the decisions}

\begin{tabularx}{\linewidth}{@{}l L p{4.4cm}@{}}
\toprule
\textbf{\#} & \textbf{Finding} & \textbf{Source} \\
\midrule
E1 & Heidi prices per housing unit per year & \href{https://heidisystems.com/blog/messdienstleister}{heidisystems.com blog} \\
E2 & Heidi integrates with Domus, Immoware24, Wodis Sigma, Immocloud & heidisystems.com blog \\
E3 & In a WEG the owners choose the metering provider by resolution, usually on the manager's
     proposal. The legal basis is \S\,19 WEG since the 2020 reform (Heidi's blog still cites
     \S\,21) & \href{https://www.buzer.de/19_WEG.htm}{\S\,19 WEG}, heidisystems.com blog \\
E4 & For rental property, the landlord decides, or the manager if authorised & heidisystems.com blog \\
E5 & Many metering contracts run for up to 10 years & heidisystems.com blog \\
E6--E7 & WBM and HOWOGE run metering through their own subsidiaries & WBM job ad (Oct 2026);
         HOWOGE tender 1333040 \\
E9 & Meters that can't be read remotely must be replaced by 31 Dec 2026. After that, monthly
     consumption information is mandatory & HeizkostenV \S\,5, \S\,6a \\
E10 & If they aren't replaced, tenants may cut their heating-cost share by 3\% & HeizkostenV \S\,12 \\
E12 & VDIV publishes no member list, so membership has to come from each company's website &
      vdiv.de/verwaltersuche \\
\bottomrule
\end{tabularx}

%--------------------------------------------------------------------
\section{Sources and tools}

\begin{tabularx}{\linewidth}{@{}p{3.6cm} L@{}}
\toprule
\textbf{Source / tool} & \textbf{Used for} \\
\midrule
Google Places API (New) & Finding companies. Berlin is split into map cells, and a cell is
  split further whenever a search hits the 60-result limit. The search itself uses free ID-only
  requests; website and phone are requested only where needed (157 billed requests, inside the
  free tier) \\
OpenStreetMap (Overpass) & A second source for companies (\code{office=property\_management}) \\
Company websites & Every stored fact comes from the company's own pages: home, Impressum,
  about, services, team, careers \\
Claude Haiku 5.5 & The research agent (Anthropic API) \\
Python & requests and tenacity (retries), BeautifulSoup (HTML), RapidFuzz (quote matching),
  pydantic (schema), openpyxl (Excel), pytest and ruff \\
\bottomrule
\end{tabularx}

%--------------------------------------------------------------------
\section{Automation: the pipeline and the agent}

The rule I followed: \textbf{the AI is used where judgement is needed, and code where the answer
must be exact.} The AI reads websites and extracts facts. It never calculates a score. One
command runs the whole pipeline with no manual steps: \code{python -m pipeline.main all}.

\begin{enumerate}
  \item \textbf{Collect} (code): Google Places and OpenStreetMap.
  \item \textbf{Merge} (code): duplicates merged by website domain, then by normalised name.
  \item \textbf{Fetch} (code): up to 6 pages per site; robots.txt respected.
  \item \textbf{Research} (AI): the agent reads the pages, may fetch up to 4 more on the same
        site, and submits each fact with its source URL and a word-for-word quote.
  \item \textbf{Verify} (code): every quote is checked against the page text. If it fails, the
        value becomes unknown.
  \item \textbf{Score and export} (code): the formula above, written to Excel.
\end{enumerate}

\subsection{What matters about the agent}
\begin{itemize}
  \item \textbf{Evidence or unknown.} Every value needs a URL and a word-for-word quote. Without
        one, the agent has to answer ``unknown''.
  \item \textbf{Checked by code, not trusted.} The code compares each quote with the page text and
        checks that numbers actually appear in the quote.
  \item \textbf{Two tools only.} \code{fetch\_page} (pages on the same company site only) and
        \code{submit\_record} (a strict schema; invalid answers are sent back once for correction).
  \item \textbf{Cheap by design.} Our code fetches the pages itself instead of using a paid web
        search, and prompt caching makes the long instructions almost free after the first
        company. The result is about \$0.002 per company.
\end{itemize}

\subsection{Guardrails}
\begin{itemize}
  \item \textbf{Budget:} a hard stop at \$4 total and \$0.05 per company, based on costs
        written to an append-only spend ledger.
  \item \textbf{Quality gate:} after 25 companies the run checks its own failure rate,
        verification rate and cost, and stops itself if any is out of range.
  \item \textbf{Restartable:} each company is saved as soon as it's done, and a restart skips it.
        One failing company never stops the run.
  \item \textbf{Polite crawling:} robots.txt is respected, requests are spaced out, and blocked
        sites are skipped, never worked around.
\end{itemize}

%--------------------------------------------------------------------
\section{Validation}

\begin{tabularx}{\linewidth}{@{}L l@{}}
\toprule
\textbf{Check} & \textbf{Result} \\
\midrule
Agent claims confirmed on the cited page & 3,945 / 4,127 (96\%) \\
Named contacts confirmed on the cited page & 1,601 / 1,686 (95\%) \\
Quality gate after 25 companies & passed: 0\% failures, 97\% verified \\
Accuracy against hand-checked companies (Haiku) & 41 / 45 fields (91\%) \\
Haiku vs Sonnet on 5 of those companies & \textcolor{assume}{pending} \\
\bottomrule
\end{tabularx}

The verification rate shows the agent doesn't invent quotes. The hand-checked test measures
something different: whether it finds the right facts and doesn't miss any. For that test I
checked 10 companies' websites by hand, without looking at the agent's answers, and the code
compares the two field by field (units, residential share, buyer type, software, association).

\begin{itemize}
  \item \textbf{One site doesn't count:} its robots.txt blocks crawlers, so the agent never saw
        it. The remaining 9 companies give 45 fields.
  \item \textbf{Units, software and association: 27 of 27 correct.} The agent also found a figure
        the hand check had missed ("more than 7,000 apartments" on an English page). The answer
        key was corrected.
  \item \textbf{The 4 misses are cautious, not invented.} Two are ``unknown'' because the agent
        didn't open the one subpage that described the services. One is ``unknown'' because its
        quote combined two list items, so the strict quote check rejected it. One is a borderline
        call (rental-focused vs.\ unclear).
  \item The agent never stated a number or membership that wasn't on the site.
\end{itemize}

%--------------------------------------------------------------------
\section{Results}

\begin{tabularx}{\linewidth}{@{}L r@{}}
\toprule
\textbf{In the sheet} & \textbf{Companies} \\
\midrule
Scored with confidence OK & 242 \\
Scored with confidence Low (size or most KPIs unknown) & 904 \\
Excluded: not a property manager (brokers, caretakers, law firms) & 417 \\
Site dead, blocked or disallowed by robots.txt & 219 \\
No website of their own, so not researched (kept with phone and address) & 517 \\
Named contact person & 1,098 \\
Named contact with a personal email & 88 \\
General company email & 1,070 \\
\bottomrule
\end{tabularx}

\subsection{What the data showed}
\begin{itemize}
  \item \textbf{Size is rarely published.} Only 178 researched companies state their units
        (median 2,000). Another 112 could be estimated from employees or buildings.
  \item \textbf{About a quarter of Google's ``property management'' results aren't
        Hausverwaltungen.} The agent marked 417 as brokers, caretakers and similar businesses,
        and they're excluded instead of scored.
  \item \textbf{Decision-maker emails are scarce.} Most sites name a managing director in the
        Impressum but only give a general office address. Personal emails turned up for 88
        contacts.
  \item \textbf{Municipal companies run their own metering.} The agent found that degewo and
        STADT UND LAND use their own subsidiaries, which supports scoring municipal companies low.
  \item \textbf{Software and metering providers are rarely mentioned.} 82 companies name their
        software, and at least 11 of them use Domus or Immoware24. Only 7 name their metering
        provider.
  \item \textbf{Associations:} 120 companies show a membership, mostly IVD and VDIV.
\end{itemize}

%--------------------------------------------------------------------
\section{Repository, structure and running}

Code: \url{\repourl}. Setup and every command are also in the repository's README.

\subsection{Project structure}
\begin{Verbatim}[fontsize=\small]
pipeline/
  main.py            one entry point; runs the steps in order
  config.py          every number: models, weights, budgets, limits
  models.py          one schema: Company, Fact (value+status+source+quote)
  collectors/        google_places.py, osm.py  (each returns list[Company])
  dedupe.py          merges duplicates; independent of input order
  fetching/          robots.txt, safe redirects, HTML to text, page cache
  agent/             research_agent.py, prompt, tools, schema, cost, ledger
  research_step.py   order, budget stop, quality gate
  checkpoint.py      saves each company; skipped on restart
  verify.py          re-checks every quote on its page
  scoring.py         the lead score (pure functions)
  evaluation.py      accuracy against hand-checked answers
  export/            Excel (Google Sheets optional)
  utils/             HTTP with retries, URL normalising, file I/O, logging
tests/               247 pytest tests
\end{Verbatim}

\subsection{Running}
\begin{Verbatim}[fontsize=\small]
py -m venv .venv
.venv\Scripts\pip install -r requirements.txt
copy .env.example .env          (then add the API keys)

.venv\Scripts\python -m pipeline.main all        whole pipeline
.venv\Scripts\python -m pipeline.main research --limit 50
.venv\Scripts\python -m pipeline.main export     rebuild the Excel file
.venv\Scripts\python -m pipeline.main evaluate   accuracy test
.venv\Scripts\python -m pytest                   tests
\end{Verbatim}

%--------------------------------------------------------------------
\section{Limits and open assumptions}
\begin{itemize}
  \item Residential share, the cooperative and private-landlord scores, and both bonus signals
        are assumptions. They should be checked against Heidi's real win data.
  \item LinkedIn wasn't searched. It needs search tools or accounts and would break LinkedIn's
        terms, so contacts come only from company websites.
  \item Google Places was used to find companies. Stored facts come from the companies' own
        websites, which fits Google's storage rules.
  \item Companies without their own website stay in the list unresearched. They could be covered
        next with directory pages or a phone call.
\end{itemize}

\begin{thebibliography}{9}
\bibitem{gqm} V.\,R.~Basili, G.~Caldiera, H.\,D.~Rombach. \emph{The Goal Question Metric
  Approach}. Encyclopedia of Software Engineering, Wiley, 1994.
\bibitem{smart} G.\,T.~Doran. \emph{There's a S.M.A.R.T. way to write management's goals and
  objectives}. Management Review 70(11), 1981.
\end{thebibliography}

\end{document}
